\chapter{Zaklju\v cak}
 Dokazi nultog znanja predstavljaju revolucionarno otkri\' ce u oblasti kriptografije, obezbeđuju\' ci sredstvo za dokazivanje znanja bez otkrivanja samog znanja. U ovom radu prikazan je kriptosistem koji predstavlja preteču dokaza nultog znanja, navedena je prva definicija ovog koncepta i osnovna podela na interaktivne i neinteraktivne dokaze nultog znanja. Na problemu kvadratnog ostatka prikazani su primeri oba tipa dokaza. Upoređena su tri najpoznatija neinteraktivna dokaza nultog znanja: \textenglish{ZK-SNARK, ZK-STARK} i \textenglish{Bulletproofs}. Pored toga izloženo je da je primena dokaza nultog znanja \v siroko rasprostranjena.
 
  Kako tehnologija nastavlja da se razvija, primena dokaza nultog znanja se sve vi\v se \v siri. Dokazi nultog znanja su danas jedan od najpopularnijih metoda za osiguravanje anonimnosti transakcija pa igraju ključnu ulogu u razvoju bezbednih digitalnih sistema koji čuvaju privatnost korisnika.

%U ovom radu prikazano je da je primena {\it\textenglish{manifold learning}}-a \v siroko rasprostranjena. Koristi se u raznim poljima poput bioinformatike, obrade prirodnog jezika i modelovanja mre\v za. Prikazani su algoritmi: Multidimenzionalno skaliranje \textenglish{MDS}, {\it\textenglish{Isomap embedding}}, Lokalno linearno potapanje \textenglish{LLE}, Laplasovo karakteristi\v cno preslikavanje i  \textenglish{t-SNE}.  Prikazane su tri alternative \textenglish{LLE} algoritma: Modifikovano lokalno linearno potapanje, Hesijanovo karakteristi\v cno preslikavanje i {\it\textenglish{Local Tangent Space Alignment }}. Od svih ovih algoritama samo su  {\it\textenglish{Isomap}}, \textenglish{MLLE}, \textenglish{HLLE} i \textenglish{LTSA} uspeli da otkriju dvodimenzionalnu strukturu mnogostrukosti koja je u radu posmatrana. \textenglish{MDS} u tome nije uspeo jer je njegova najve\' ca slabost no\v senje sa nelinearnim kompleksnim strukturama. Iz sli\v cnih razloga i ostali algoritmi nisu uspeli da uhvate globalnu strukturu mnogostrukosti. Pored toga jako je bitno dobro podesiti parametre da bi se na izlazu iz algoritma dobio skup podataka koji o\v cuvava sve \v zeljene osobine po\v cetnih podataka.